\documentclass[10pt]{extarticle}
\usepackage[a4paper,margin=1.5cm]{geometry}
\usepackage[T1]{fontenc}
\usepackage{amsmath,amssymb}
\usepackage{graphicx}
\usepackage[export]{adjustbox}
\usepackage{multicol}
\usepackage{enumitem}
\usepackage{booktabs,array,tabularx}
\setlength{\columnsep}{0.7cm}
\setlength{\parindent}{0pt}
\setlength{\parskip}{3pt}
\pagestyle{empty}
\makeatletter
\newcommand{\qfigdef}[4]{\@namedef{qf@file@#1}{#2}\@namedef{qf@w@#1}{#3}\@namedef{qf@wide@#1}{#4}}
\newcommand{\qf@check}[1]{\@ifundefined{qf@file@#1}{\PackageError{score2}{Unknown figure id #1}{}}{}}
\newcommand{\qf@img}[4][]{\includegraphics[width=#2\dimexpr\@nameuse{qf@w@#4}pt\relax,max width=#3,max height=0.45\textheight,#1]{\@nameuse{qf@file@#4}}}
\newif\ifqf@full
\newif\ifqf@boxed
\newsavebox{\qf@box}
\newcommand{\qf@open}{\global\qf@fulltrue\par\global\qf@boxedtrue
  \setbox\qf@box\vbox\bgroup\hsize\textwidth\linewidth\textwidth}
\newcommand{\qf@emit}[1]{\par\egroup\global\qf@boxedfalse
  \ifdim\dimexpr\ht\qf@box+\dp\qf@box\relax<#1\textheight\noindent\usebox{\qf@box}\par\else\unvbox\qf@box\fi}
\newcommand{\qfullsplit}{\ifqf@boxed\qf@emit{0.85}\fi}
\newcommand{\qf@close}{\ifqf@boxed\qf@emit{0.45}\fi\global\qf@fullfalse}
\newcommand{\qfullbegin}{\end{multicols}\qf@open}
\newcommand{\qfullend}{\qf@close\begin{multicols}{2}}
\newcommand{\qfullstart}{\qf@open}
\newcommand{\qfulldone}{\qf@close}
\newcommand{\qfig}[2][1]{\qf@check{#2}%
  \ifqf@full\qf@blk{#1}{#2}\else
  \ifnum\@nameuse{qf@wide@#2}=1 \expandafter\@firstoftwo\else\expandafter\@secondoftwo\fi
  {\end{multicols}\begin{center}\qf@img{#1}{\textwidth}{#2}\end{center}\begin{multicols}{2}}%
  {\qf@blk{#1}{#2}}\fi}
\newcommand{\qf@blk}[2]{\par\smallskip{\centering\qf@img{#1}{\linewidth}{#2}\par}\smallskip}
\newcommand{\qfigi}[2][1]{\qf@check{#2}\qf@img[valign=c]{#1}{\linewidth}{#2}}
\makeatother
\newcommand{\papersec}[1]{\par\bigskip{\centering\Large\bfseries PAPER-#1\par}\medskip}
\newcommand{\partsec}[2]{\par\medskip\noindent{\large\bfseries PART-#1: #2}\par\smallskip}
\newcommand{\sect}[1]{\par\smallskip\noindent\textbf{#1}\par}
\newcommand{\qn}[1]{\par\smallskip\noindent\textbf{#1.}~}
\newcommand{\paraq}[2]{\par\smallskip\noindent\textbf{Paragraph for Questions #1 and #2}\par}
\newenvironment{qoptions}{\begin{itemize}[leftmargin=1.8em,itemsep=1pt,topsep=2pt,label={}]}{\end{itemize}}
\newcommand{\opt}[2]{\item[(#1)] #2}
% figure registry (generated by score2_tools.py finalize)
\qfigdef{p1_phy_s1_q01_1}{Score-I - Test 05 (22 Jan 2022)-img/p1_phy_s1_q01_1.png}{88.0}{0}
\qfigdef{p1_phy_s1_q06_1}{Score-I - Test 05 (22 Jan 2022)-img/p1_phy_s1_q06_1.png}{133.0}{0}
\qfigdef{p1_phy_s2_q07_1}{Score-I - Test 05 (22 Jan 2022)-img/p1_phy_s2_q07_1.png}{44.9}{0}
\qfigdef{p1_phy_s2_q08_1}{Score-I - Test 05 (22 Jan 2022)-img/p1_phy_s2_q08_1.png}{174.4}{0}
\qfigdef{p1_phy_s2_q09_1}{Score-I - Test 05 (22 Jan 2022)-img/p1_phy_s2_q09_1.png}{154.5}{0}
\qfigdef{p1_phy_s2_q10_1}{Score-I - Test 05 (22 Jan 2022)-img/p1_phy_s2_q10_1.png}{139.4}{0}
\qfigdef{p1_phy_s3_q01_1}{Score-I - Test 05 (22 Jan 2022)-img/p1_phy_s3_q01_1.png}{151.3}{0}
\qfigdef{p1_phy_s3_q02_1}{Score-I - Test 05 (22 Jan 2022)-img/p1_phy_s3_q02_1.png}{146.7}{0}
\qfigdef{p1_phy_s3_q03_1}{Score-I - Test 05 (22 Jan 2022)-img/p1_phy_s3_q03_1.png}{124.6}{0}
\qfigdef{p1_phy_s3_q04_1}{Score-I - Test 05 (22 Jan 2022)-img/p1_phy_s3_q04_1.png}{128.8}{0}
\qfigdef{p1_phy_s3_q05_1}{Score-I - Test 05 (22 Jan 2022)-img/p1_phy_s3_q05_1.png}{117.6}{0}
\qfigdef{p1_phy_s3_q07_1}{Score-I - Test 05 (22 Jan 2022)-img/p1_phy_s3_q07_1.png}{122.4}{0}
\qfigdef{p1_phy_s3_q08_1}{Score-I - Test 05 (22 Jan 2022)-img/p1_phy_s3_q08_1.png}{82.0}{0}
\qfigdef{p1_chem_s1_q02_A}{Score-I - Test 05 (22 Jan 2022)-img/p1_chem_s1_q02_A.png}{32.5}{0}
\qfigdef{p1_chem_s1_q02_B}{Score-I - Test 05 (22 Jan 2022)-img/p1_chem_s1_q02_B.png}{38.9}{0}
\qfigdef{p1_chem_s1_q02_C}{Score-I - Test 05 (22 Jan 2022)-img/p1_chem_s1_q02_C.png}{60.5}{0}
\qfigdef{p1_chem_s1_q02_D}{Score-I - Test 05 (22 Jan 2022)-img/p1_chem_s1_q02_D.png}{54.6}{0}
\qfigdef{p1_chem_s2_q10_P}{Score-I - Test 05 (22 Jan 2022)-img/p1_chem_s2_q10_P.png}{32.6}{0}
\qfigdef{p1_chem_s2_q10_Q}{Score-I - Test 05 (22 Jan 2022)-img/p1_chem_s2_q10_Q.png}{63.9}{0}
\qfigdef{p1_chem_s2_q10_R}{Score-I - Test 05 (22 Jan 2022)-img/p1_chem_s2_q10_R.png}{49.7}{0}
\qfigdef{p1_chem_s2_q10_S}{Score-I - Test 05 (22 Jan 2022)-img/p1_chem_s2_q10_S.png}{56.6}{0}
\qfigdef{p1_chem_s3_q03_1}{Score-I - Test 05 (22 Jan 2022)-img/p1_chem_s3_q03_1.png}{264.0}{1}
\qfigdef{p1_chem_s3_q05_i}{Score-I - Test 05 (22 Jan 2022)-img/p1_chem_s3_q05_i.png}{64.1}{0}
\qfigdef{p1_chem_s3_q05_ii}{Score-I - Test 05 (22 Jan 2022)-img/p1_chem_s3_q05_ii.png}{78.2}{0}
\qfigdef{p1_chem_s3_q05_iii}{Score-I - Test 05 (22 Jan 2022)-img/p1_chem_s3_q05_iii.png}{31.1}{0}
\qfigdef{p1_chem_s3_q05_iv}{Score-I - Test 05 (22 Jan 2022)-img/p1_chem_s3_q05_iv.png}{54.2}{0}
\qfigdef{p1_chem_s3_q05_v}{Score-I - Test 05 (22 Jan 2022)-img/p1_chem_s3_q05_v.png}{54.4}{0}
\qfigdef{p1_chem_s3_q05_vi}{Score-I - Test 05 (22 Jan 2022)-img/p1_chem_s3_q05_vi.png}{67.6}{0}
\qfigdef{p1_chem_s3_q05_vii}{Score-I - Test 05 (22 Jan 2022)-img/p1_chem_s3_q05_vii.png}{63.5}{0}
\qfigdef{p1_chem_s3_q05_viii}{Score-I - Test 05 (22 Jan 2022)-img/p1_chem_s3_q05_viii.png}{45.5}{0}
\qfigdef{p1_chem_s3_q05_ix}{Score-I - Test 05 (22 Jan 2022)-img/p1_chem_s3_q05_ix.png}{43.0}{0}
\qfigdef{p1_chem_s3_q05_x}{Score-I - Test 05 (22 Jan 2022)-img/p1_chem_s3_q05_x.png}{33.9}{0}
\qfigdef{p1_math_s2_q07_S}{Score-I - Test 05 (22 Jan 2022)-img/p1_math_s2_q07_S.png}{80.3}{0}

\begin{document}
\begin{center}{\Large\bfseries JEE (Main+Advanced) Score-I --- Test 05 (22 Jan 2022)}\end{center}
\begin{multicols}{2}
\raggedcolumns
\partsec{1}{PHYSICS}
\sect{Section-I(i)}
\qn{1} A square plate of mass `M' and edge `a' is hinged at one of its vertex and is allowed to fall down in a vertical plane as shown in figure:
\qfig{p1_phy_s1_q01_1}
\begin{qoptions}
\opt{A}{The reaction along the line joining centre \& hinge becomes zero when the plate moves by an angle $\theta=53^\circ$ from vertically upward line.}
\opt{B}{The hinge reaction when the centre comes into the same horizontal line as hinge is $\dfrac{mg}{4}$}
\opt{C}{The max angular acceleration of the plate is $\dfrac{3g}{2\sqrt2a}$}
\opt{D}{The max speed of the COM of the plate is $\sqrt{1.5ag}$}
\end{qoptions}
\qn{2} A solid sphere of mass M and radius R is allowed to roll down a rough inclined plane of angle $\theta$. The centre of the sphere is at a height `h' above the ground.
\begin{qoptions}
\opt{A}{If $\mu=\dfrac27\tan\theta$, the solid sphere can roll without slipping.}
\opt{B}{If $\mu=\dfrac{\tan\theta}7$, the time taken by sphere to reach the bottom is $\dfrac1{\sin\theta}\sqrt{\dfrac{7h}{3g}}$}
\opt{C}{If $\mu=\dfrac{\tan\theta}7$, the angular speed of the sphere at the bottom is $\sqrt{\dfrac{10g}{7R}}$}
\opt{D}{If $\mu=\dfrac{\tan\theta}7$ the angular acceleration of the sphere is $\dfrac{5g\sin\theta}{14R}$}
\end{qoptions}
\qn{3} A disk of mass M and radius `R' is moving on a rough horizontal surface with COM speed $v_0$ and some angular speed. If P is the point of contact of the disk with ground and Q is the highest point then
\begin{qoptions}
\opt{A}{for $\omega_0=\dfrac{v_0}R$ the magnitude of angular momentum about point P would always be $\dfrac32mv_0R$}
\opt{B}{for $\omega_0=\dfrac{v_0}R$ the magnitude of angular momentum about point `P' can be $\dfrac{mv_0R}2$}
\opt{C}{for $\omega_0=\dfrac{v_0}R$ the angular momentum about a point taken on ground at a distance more than `2R' from centre should remain conserved}
\opt{D}{for $\omega_0=\dfrac{v_0}R$ the disk should roll without slipping.}
\end{qoptions}
\qn{4} The speed of light in glass lens is $2\times10^8$ ms$^{-1}$. If a plano convex lens of apprature 10 cm and thickness at the middle ($t=2$mm) is taken then
\begin{qoptions}
\opt{A}{the refractive index of the lens is 1.5}
\opt{B}{the radius of curvature of the curved surface is nearly 125 cm}
\opt{C}{the focal length of the lens is nearly 125 cm}
\opt{D}{all of the above}
\end{qoptions}
\qn{5} A plano convex of refractive index $\dfrac32$ is silvered on plane surface and it behaves like a concave mirror of focal length 30cm, then
\begin{qoptions}
\opt{A}{the radius of curvature of the curved surface is 30 cm}
\opt{B}{If the curved surface is silvered the focal length would become 10cm}
\opt{C}{If the lens is dipped in water it would behave like a diverging lens.}
\opt{D}{The lens can behave like a converging lens even when the medium around it changes.}
\end{qoptions}
\qn{6} Two plane mirrors M$_1$ and M$_2$ are inclined at an angle `$\theta$' which can be varied by rotating the mirrors as shown in figure.
\qfig{p1_phy_s1_q06_1}
\begin{qoptions}
\opt{A}{If a ray coming parallel to M$_2$ first strikes M$_1$ then M$_2$ and then only after striking M$_1$ retraces its complete path then $\theta=30^\circ$.}
\opt{B}{If a ray of light strikes the mirror M$_1$ at $30^\circ$ from the normal to M$_1$ then angle of deviation suffered by ray is always $360^\circ$ (consider the rays strikes only once on both mirrors)}
\opt{C}{If angle between the mirror is $30^\circ$ and the ray strikes mirror M$_1$ at $70^\circ$ to the mirror then net deviation suffered by the ray is $300^\circ$. (consider the rays strikes only once on both mirror).}
\opt{D}{None of the above}
\end{qoptions}
\sect{Section-I(ii)}
\qn{7} A uniform rod of length L is moving with speed $v$ on a horizontal smooth surface. One end of the rod is suddenly hinged - this process takes a very small time $L/100v$. Match the ratios in list I with their values in list II.
\qfig{p1_phy_s2_q07_1}
\begin{center}\small\setlength{\tabcolsep}{3pt}
\begin{tabular}{|c|p{4.2cm}|c|c|}
\hline
\multicolumn{2}{|c|}{List-I} & \multicolumn{2}{c|}{List-II}\\ \hline
(P) & Speed of the centre of mass of the rod just after to just before the event. & (1) & 2/3\\ \hline
(Q) & KE of the rod just after to just before & (2) & 2/9\\ \hline
(R) & Average force exerted by the hinge during hinging to (100 $\times$ force exerted by hinge after hinging) & (3) & 3/4\\ \hline
(S) & Time taken to turn by the rod one radian after the event to time taken to cover the length of the rod before the event. & (4) & 5/6\\ \hline
\end{tabular}
\end{center}
\begin{qoptions}
\opt{A}{P $\to$ 4;Q $\to$ 2;R $\to$ 1;S $\to$ 3}
\opt{B}{P $\to$ 3;Q $\to$ 2;R $\to$ 4;S $\to$ 1}
\opt{C}{P $\to$ 3;Q $\to$ 3;R $\to$ 2;S $\to$ 1}
\opt{D}{P $\to$ 1;Q $\to$ 4;R $\to$ 3;S $\to$ 2}
\end{qoptions}
\qn{8} An object lies on the principal axis of a convex mirror of focal length $f=1.0$ m at a distance $3f$ from it. At time $t=0$, the object starts moving with velocity components $v_0=1$ m/s along the principal axis (towards the mirror) and $v_0=1$ m/s perpendicular to the principal axis. The image moves along the line $ax+by=c$ (where a, b and c are relatively prime integers and $x$ and $y$ are in meters.) Match the quantities given in List--I with their magnitudes (in SI units) in List--II. We are only concerned with the time in which the rays from the object to the mirror are paraxial.
\qfig{p1_phy_s2_q08_1}
\begin{center}\small\setlength{\tabcolsep}{3pt}
\begin{tabular}{|c|p{4.8cm}|c|c|}
\hline
\multicolumn{2}{|c|}{List-I} & \multicolumn{2}{c|}{List-II}\\ \hline
(P) & Axial component of velocity of image at time $t=0$ & (1) & 1/3\\ \hline
(Q) & Transverse component of velocity of image at time $t=0$ & (2) & 1/16\\ \hline
(R) & The value of $1/a$ & (3) & 1/4\\ \hline
(S) & The value of $1/c$ & (4) & 1/6\\ \hline
\end{tabular}
\end{center}
\begin{qoptions}
\opt{A}{P $\to$ 2;Q $\to$ 3;R $\to$ 3;S $\to$ 1}
\opt{B}{P $\to$ 1;Q $\to$ 3;R $\to$ 4;S $\to$ 2}
\opt{C}{P $\to$ 2;Q $\to$ 4;R $\to$ 4;S $\to$ 1}
\opt{D}{P $\to$ 1;Q $\to$ 3;R $\to$ 2;S $\to$ 4}
\end{qoptions}
\qn{9} A uniform disc of radius $R=0.80$ m rolls toward a vertical step of height $R/6$ with speed $u$. When it hits the step, it does not bounce or slip. In order for it to climb the step, we need $u>u_1$ and for the disc not to bounce, we need $u<u_2$. Taking $g=10$ m/s$^2$ and $\sqrt{15}=3.88$ match the quantities mentioned in List--I with their magnitudes in SI units in List--II.
\qfig{p1_phy_s2_q09_1}
\begin{center}\small\setlength{\tabcolsep}{3pt}
\begin{tabular}{|c|p{4.8cm}|c|c|}
\hline
\multicolumn{2}{|c|}{List-I} & \multicolumn{2}{c|}{List-II}\\ \hline
(P) & If $u=1$ m/s, the angular velocity of the disc immediately after impact & (1) & 7/27\\ \hline
(Q) & $u_1$ & (2) & 1.41\\ \hline
(R) & $u_2-u_1$ & (3) & 1.50\\ \hline
(S) & If $u=1$ m/s and mass of disc $=$ 1 kg, the horizontal component of the impulse that the disc experiences from the step. & (4) & 10/9\\ \hline
\end{tabular}
\end{center}
\begin{qoptions}
\opt{A}{P $\to$ 3;Q $\to$ 1;R $\to$ 4;S $\to$ 2}
\opt{B}{P $\to$ 4;Q $\to$ 3;R $\to$ 2;S $\to$ 1}
\opt{C}{P $\to$ 2;Q $\to$ 4;R $\to$ 1;S $\to$ 3}
\opt{D}{P $\to$ 1;Q $\to$ 2;R $\to$ 3;S $\to$ 4}
\end{qoptions}
\qn{10} A slab occupies the region between $y=0$ and $y=R=1$ m. Rays of light traveling in the XY plane enter the slab at the origin at angle of incidence $i$ and are found to travel in circular paths. Consider the following cases:

Case 1 : The refractive index of the slab, $\mu$, varies from 1 at $y=0$ to infinity at $y=R$ and $i=90$.

Case 2 : The refractive index of the slab, $\mu$, varies from infinity at $y=0$ to 1 at $y=R$ and $i=0$.
\qfig{p1_phy_s2_q10_1}
\begin{center}\small\setlength{\tabcolsep}{3pt}
\begin{tabular}{|c|p{3.4cm}|c|c|}
\hline
\multicolumn{2}{|c|}{List-I} & \multicolumn{2}{c|}{List-II}\\ \hline
(P) & In case 1, $1/\mu$ at $y=R/4$ & (1) & 4/5\\ \hline
(Q) & In case 1, $x$ at $y=2R/5$ & (2) & 1/5\\ \hline
(R) & In case 2, $1/\mu$ at $y=R/2$ & (3) & 3/4\\ \hline
(S) & In case 2, $x$ at $y=3R/5$ & (4) & 1/2\\ \hline
\end{tabular}
\end{center}
\begin{qoptions}
\opt{A}{P $\to$ 4;Q $\to$ 1;R $\to$ 2;S $\to$ 3}
\opt{B}{P $\to$ 1;Q $\to$ 4;R $\to$ 3;S $\to$ 2}
\opt{C}{P $\to$ 2;Q $\to$ 2;R $\to$ 1;S $\to$ 4}
\opt{D}{P $\to$ 3;Q $\to$ 1;R $\to$ 4;S $\to$ 2}
\end{qoptions}
\sect{Section-II}
\qn{1} Isosceles triangle glass prism stands with its base in water. The angle that its two equal sides makes with base are $\theta$ each. An incident ray of light parallel to the water surface internally reflect at the glass water interface and subsequently re-emerges into the air, for this $\theta$ must be at least $\tan^{-1}x$. Find the value of $x$. ($\mu_g=\dfrac53$, $\mu_w=\dfrac43$)
\qfig{p1_phy_s3_q01_1}
\qn{2} A light ray I is incident on a plane mirror M. The mirror is rotated in the direction as shown in the figure by an arrow at frequency $\dfrac2\pi$ rps. The light reflected by the mirror is received on the wall W at a distance 5 m from the axis of rotation. When the angle of incidence become $37^\circ$ the speed (in m/s) of a spot on the wall is
\qfig{p1_phy_s3_q02_1}
\qn{3} Centre of a circular disc of uniform density of radius $R=0.5$m and mass $M=81$ kg is at O. The disc may be assumed to have two part one is another circular disc C of radius $R/3$ with centre at $O_1$ $(OO_1=2R/3)$ and the other part K, has its centre of mass at $O_2$. Moment of inertia of disc K about an axis perpendicular to this plane of the disc and passing through $O_2$ is
\qfig{p1_phy_s3_q03_1}
\qn{4} Rod is released from the position as shown. Find hinge force (in N) when it becomes vertical.
\qfig{p1_phy_s3_q04_1}
\qn{5} A disc of Mass `M' $=6$kg and Radius $R=3$m is lying on a horizontal smooth surface free to move. A particle of mass $m$ moving with velocity $V_0=5.2$ m/s strikes at the surface tangentially as shown in fig. Calculate velocity (in m/s) of centre of disc (for $e=0$)
\qfig{p1_phy_s3_q05_1}
\qn{6} A uniform thin rod of mass $M=2$ kg and length $\ell=2$m is standing on a smooth horizontal surface. A slight disturbance cause the lower end to slip on the smooth surface and the rod start falling. Find the velocity of centre of mass (in m/s) of the rod at the instant when rod makes an angle $53^\circ$ with horizontal.
\qn{7} In the given figure `O' is a point object kept on the principal axis of a concave mirror of radius of curvature 20 cm, prism of angle $1.8^\circ$. Light falling on the prism (at a small angle of incidence) get refracted through prism and then fall on the mirror. Refractive index of prism is 3/2. Find the distance (in cm) between image formed by the concave mirror due to this light.
\qfig{p1_phy_s3_q07_1}
\qn{8} A uniform rod of length 0.72 m is pivoted at one end and is hanging. It receive a horizontal impulse at a distance $x$ below the pivot as shown. If the impulse applied by the pivot is zero then find the value of $x$ (in m).
\qfig{p1_phy_s3_q08_1}
\raggedcolumns
\partsec{2}{CHEMISTRY}
\sect{Section-I(i)}
\qn{1} (A) : $\mathrm{CH_3O-CH{=}CH-CH_2Br}$ \qquad (B) : $\begin{array}[t]{@{}l@{}}\ \ \ \mathrm{OCH_3}\\[-2pt]\ \ \ \ \ |\\[-2pt]\mathrm{CH_2{=}C-CH_2Br}\end{array}$

Which of the following is(are) true about compounds (A) and (B)?
\begin{qoptions}
\opt{A}{Solvolysis of compound (A) is faster than compound (B).}
\opt{B}{S$_\mathrm{N}$2 reaction of compound (A) is faster than compound (B).}
\opt{C}{S$_\mathrm{N}$2 reaction of compound (B) is faster than compound (A).}
\opt{D}{Compound (A) undergoes reaction under unimolecular as well as bimolecular conditions.}
\end{qoptions}
\qn{2} Which of the following reactions are correctly matched with respective mechanisms?
\begin{qoptions}
\opt{A}{\qfigi{p1_chem_s1_q02_A}$+\ \mathrm{EtO^-Na^+}\longrightarrow\mathrm{E_2}$}
\opt{B}{\qfigi{p1_chem_s1_q02_B}$+\ \mathrm{NaN_3}\xrightarrow{\text{DMF}}\mathrm{S_N2}$}
\opt{C}{\qfigi{p1_chem_s1_q02_C}$+\ \mathrm{CH_3O^-Na^+}\longrightarrow\mathrm{S_N2}$}
\opt{D}{\qfigi{p1_chem_s1_q02_D}$+\ \mathrm{CH_3OH}\longrightarrow\mathrm{S_N1}$}
\end{qoptions}
\qn{3} Unknown salt (M) $+$ conc. H$_2$SO$_4$ $+$ K$_2$Cr$_2$O$_7$ $\xrightarrow[\text{Heated}]{\text{Gently}}$ Reddish brown vapour (X)

\hspace{2em}$\downarrow$ Passed into NaOH solution

Yellow ppt. (Z) $\xleftarrow[\mathrm{CH_3COOH}]{\mathrm{Pb(OAc)_2}}$ Yellow solution (Y)

The salt M is/are
\begin{qoptions}
\opt{A}{KCl}
\opt{B}{SnCl$_2$}
\opt{C}{SnCl$_4$}
\opt{D}{NaCl}
\end{qoptions}
\qn{4} Which of the following statements is/are true ?
\begin{qoptions}
\opt{A}{Fe$^{3+}$ gives red precipitate with dimethyl glyoxime in alkaline solution.}
\opt{B}{Cu$^{2+}$ ions produce a white precipitate when KCN solution is added in a small quantity.}
\opt{C}{Hg$^{2+}$ ions give deep blue precipitate with cobalt acetate and ammonium thiocyanate.}
\opt{D}{Ag$_2$O dissolve in nitric acid and also in ammonia solution.}
\end{qoptions}
\qn{5} 16.2 g of an ionic solid (100\% dissociated in water) is dissolved in 1L water. The osmotic pressure at 27$^\circ$ C for the solution is 60 R(R is universal gas constant ). Assume the density of cubic unit cell of ionic solid is 40 g/cm$^3$ and anion form the lattice in which octahedral void at body centre is vacant and cell is electrically neutral. If edge length of lattice of ionic solid is 3 \AA\ and molar mass 162 g/mol, select incorrect statement(s) (Use $N_A=6\times10^{23}$).
\begin{qoptions}
\opt{A}{100\% tetrahedral voids are occupied}
\opt{B}{100\% octahedral voids are occupied}
\opt{C}{50\% tetrahedral voids are occupied}
\opt{D}{50\% octahedral voids are occupied}
\end{qoptions}
\qn{6} Which of the following statement(s) is/are correct : (Assume ideal behaviour and identical conditions)
\begin{qoptions}
\opt{A}{Rate of effusion is directly proportional to area of hole through which gas is escaping}
\opt{B}{Helium escapes at a rate 2.65 times as fast as CO does.}
\opt{C}{Helium escapes at a rate 4 times as fast as CO$_2$ does.}
\opt{D}{Helium escapes at a rate 4 times as fast as SO$_2$ does.}
\end{qoptions}
\sect{Section-I(ii)}
\qn{7} Consider a 1.0-L container of neon gas initially at STP. Will the average kinetic energy and frequency of collisions of gas molecules with the walls of the container increase, decrease, or remain the same under each of the following conditions?
\begin{center}\small\setlength{\tabcolsep}{3pt}
\begin{tabular}{|c|p{3.4cm}|c|p{2.5cm}|}
\hline
\multicolumn{2}{|c|}{List-I} & \multicolumn{2}{c|}{List-II}\\ \hline
(P) & The temperature is increased to 100$^\circ$C. (at constant V) & (1) & Average Kinetic Energy Per mol Increase\\ \hline
(Q) & The temperature is decreased to $-50^\circ$C. (at constant V) & (2) & Same Average Kinetic Energy Per mol\\ \hline
(R) & The volume is decreased to 0.5 L. (at constant T) & (3) & wall collision frequency increase\\ \hline
(S) & The number of moles of neon is doubled. (at constant T and V) & (4) & wall collision frequency decrease\\ \hline
 & & (5) & Average Kinetic Energy Per mol decrease\\ \hline
\end{tabular}
\end{center}
\begin{qoptions}
\opt{A}{P $\to$ 1,3;Q $\to$ 4,5;R $\to$ 2,3;S $\to$ 2, 3}
\opt{B}{P $\to$ 2,3;Q $\to$ 4,5;R $\to$ 3,4;S $\to$ 1,3}
\opt{C}{P $\to$ 1,4;Q $\to$ 2,3;R $\to$ 4,5;S $\to$ 3,5}
\opt{D}{P $\to$ 1,3;Q $\to$ 1,4;R $\to$ 1,4;S $\to$ 2,5}
\end{qoptions}
\qn{8} Van der Waals constant a and b for a gas are 4.0 atm L$^2$/mol$^2$ and 0.05 L/mol, respectively. If 80g of gas (molecular mass $=16$) is placed in a 10 L vessel at 300 K, then
\begin{center}\small\setlength{\tabcolsep}{3pt}
\begin{tabular}{|c|p{3.4cm}|c|c|}
\hline
\multicolumn{2}{|c|}{List-I} & \multicolumn{2}{c|}{List-II}\\ \hline
(P) & Pressure correction (in atm) & (1) & 0.0625\\ \hline
(Q) & Free volume of gas (in L) & (2) & 1\\ \hline
(R) & Actual volume occupied by molecules (in L) & (3) & 9.75\\ \hline
(S) & Volume correction (in L) & (4) & 0.25\\ \hline
\end{tabular}
\end{center}
\begin{qoptions}
\opt{A}{P $\to$ 1;Q $\to$ 2;R $\to$ 3;S $\to$ 4}
\opt{B}{P $\to$ 2;Q $\to$ 3;R $\to$ 4;S $\to$ 1}
\opt{C}{P $\to$ 2;Q $\to$ 3;R $\to$ 1;S $\to$ 4}
\opt{D}{P $\to$ 3;Q $\to$ 4;R $\to$ 1;S $\to$ 3}
\end{qoptions}
\qn{9} Borax Bead test is carried out when the original mixture is coloured. It is done with the help of a clean platinum wire on which a small loop is made at the end. When borax is heated on platinum wire loop a transparent glass like bead is obtained. The hot bead is brought in contact with salt till it reacts with fused borax and colour is imparted to the bead. Bead colour is noted.
\begin{center}\small\setlength{\tabcolsep}{3pt}
\begin{tabular}{|c|c|c|c|}
\hline
\multicolumn{2}{|c|}{List-I} & \multicolumn{2}{c|}{List-II}\\ \hline
(P) & Blue Green & (1) & Mn$^{2+}$\\ \hline
(Q) & Yellow & (2) & Cr$^{3+}$\\ \hline
(R) & Green & (3) & Fe$^{3+}$\\ \hline
(S) & Violet & (4) & Cu$^{2+}$\\ \hline
\end{tabular}
\end{center}
\begin{qoptions}
\opt{A}{P $\to$ 1;Q $\to$ 3;R $\to$ 2;S $\to$ 4}
\opt{B}{P $\to$ 4;Q $\to$ 3;R $\to$ 2;S $\to$ 1}
\opt{C}{P $\to$ 4;Q $\to$ 1;R $\to$ 3;S $\to$ 2}
\opt{D}{P $\to$ 1;Q $\to$ 4;R $\to$ 3;S $\to$ 3}
\end{qoptions}
\qn{10} Match the following reaction in List I with appropriate reaction in List II and select the correct answer using the code given below the List:
\begin{center}\small\setlength{\tabcolsep}{3pt}
\begin{tabular}{|c|p{4.6cm}|c|c|}
\hline
\multicolumn{2}{|c|}{List-I} & \multicolumn{2}{c|}{List-II}\\ \hline
(P) & \qfigi{p1_chem_s2_q10_P}$\xrightarrow[\Delta]{\text{NaOEt}}$ & (1) & S$_\mathrm{N}$2\\ \hline
(Q) & \qfigi{p1_chem_s2_q10_Q}$\xrightarrow{\mathrm{SH^-}}$ & (2) & E$_2$\\ \hline
(R) & \qfigi{p1_chem_s2_q10_R}$\xrightarrow[\Delta]{\mathrm{Me_3C-OH}}$ & (3) & S$_\mathrm{N}$1\\ \hline
(S) & \qfigi{p1_chem_s2_q10_S}$\xrightarrow{\text{MeOH}}$ & (4) & E$_1$\\ \hline
\end{tabular}
\end{center}
\begin{qoptions}
\opt{A}{P $\to$ 2;Q $\to$ 1;R $\to$ 4;S $\to$ 3}
\opt{B}{P $\to$ 2;Q $\to$ 1;R $\to$ 3;S $\to$ 4}
\opt{C}{P $\to$ 1;Q $\to$ 3;R $\to$ 2;S $\to$ 4}
\opt{D}{P $\to$ 1;Q $\to$ 3;R $\to$ 4;S $\to$ 2}
\end{qoptions}
\sect{Section-II}
\qn{1} A solid `AB' has rock-salt structure. The radius ratio $\left(\dfrac{r_A^+}{r_B^-}\right)$ is 0.50. The packing fraction of the solid is.
\qn{2} A 100 L flask contains a mixture of methane (CH$_4$) and argon gases at 27$^\circ$C. The mass of argon present is 260g and the mole fraction of methane in the mixture is 0.750. The total kinetic energy of the gaseous mixture is (in Kcal).(R $=2$ Cal k$^{-1}$ mol$^{-1}$)
\qfullbegin
\qn{3} How many of the following alkyl halides can undergo reaction by S$_\mathrm{N}$2 mechanism under suitable conditions?
\qfig{p1_chem_s3_q03_1}
\qfullend
\qn{4} How many optically active monobromides are possible on bromination of a mixture containing all isomers of C$_6$H$_{14}$
\qn{5} In how many of the following reactions, the rate of reaction increases by increasing the concentration of nucleophile?

(i) \qfigi{p1_chem_s3_q05_i}$+\ \mathrm{CH_3O^-}\longrightarrow$

(ii) \qfigi{p1_chem_s3_q05_ii}$+\ \mathrm{CH_3S^-}\longrightarrow$

(iii) \qfigi{p1_chem_s3_q05_iii}$+\ \mathrm{CH_3OH}\longrightarrow$

(iv) \qfigi{p1_chem_s3_q05_iv}$\xrightarrow{\mathrm{C_2H_5OH}}$

(v) \qfigi{p1_chem_s3_q05_v}$\xrightarrow{\mathrm{CH_3O^-Na^+}}$

(vi) \qfigi{p1_chem_s3_q05_vi}$\xrightarrow{\mathrm{H_2O}}$

(vii) \qfigi{p1_chem_s3_q05_vii}$\xrightarrow{\text{EtOH}}$

(viii) \qfigi{p1_chem_s3_q05_viii}$\xrightarrow[\mathrm{CH_3COCH_3}]{\text{NaI}}$

(ix) \qfigi{p1_chem_s3_q05_ix}$\xrightarrow[\mathrm{H_2O}]{\text{AgF}}$

(x) \qfigi{p1_chem_s3_q05_x}$\xrightarrow{\mathrm{CH_3S^-Na^+}}$
\qn{6} Among the following statements

(i) Ag$_2$SO$_3$ is insoluble in dilute HNO$_3$.

(ii) Nitrogen dioxide (NO$_2$) is absorbed by ferrous sulphate solution forming brown solution.

(iii) Heavy metal chlorides like AgCl, HgCl$_2$. etc. do not respond to chromyl chloride test.

(iv) The group reagent for III group basic radicals is NaOH.

(v) Green colour of alkaline tetrahydroxidochromate (III) becomes yellow on adding H$_2$O$_2$.

(vi) Titan yellow solution gives red colouration/ppt. with a neutral solution containing Mg$^{2+}$ ions

(vii) Solution of nitrite is decomposed by sulphamic acid.

(viii) Fe$^{2+}$ ions give brown colour precipitate with $[\mathrm{Fe(CN)_6}]^{3-}$ ions solution.

(ix) Green precipitate of Cr(OH)$_3$ is soluble in excess Na$_2$O$_2$.

If a is total number of correct while b is total number of incorrect statements find $\dfrac{a^2+b^2}{3}$.
\qn{7} In the following reaction sequence in aqueous solution, the molar mass (g/mole) of precipitate Z is (Given atomic mass of Ag $=107.87$ g/mol, S $=32$ g/mol, O $=16$ g/mol, H $=1$ g/mol)

{\footnotesize $\mathrm{S_2O_3^{2-}}\xrightarrow{\ \mathrm{Ag^+}\ }\underset{\text{Clear Solution}}{\mathrm X}\xrightarrow{\ \mathrm{Ag^+}\ }\underset{\text{White Precipitate}}{\mathrm Y}\xrightarrow{\ \text{with time}\ }\underset{\text{black precipitate}}{\mathrm Z}$}
\qn{8} Among PbS, CdS, As$_2$S$_3$, CuS, HgS, Ag$_2$S, NiS, CoS, Bi$_2$S$_3$, SnS$_2$, ZnS

$a$ is number of black colour sulphide

$b$ is number of white colour sulphide

$c$ is number of yellow colour sulphide

Find the value of $\dfrac{a^2-b}{c}$.
\raggedcolumns
\partsec{3}{MATHEMATICS}
\sect{Section-I(i)}
\qn{1} Let $\vec{A}=\hat{i}-2\hat{j}+3\hat{k}$, $\vec{B}=2\hat{i}+\hat{j}-\hat{k}$, $\vec{C}=\hat{j}+\hat{k}$. If the vector $\vec{B}\times\vec{C}$ can be expressed as $\vec{B}\times\vec{C}=x\vec{A}+y\vec{B}+z\vec{C}$ where $x,y,z$ are scalars, then which of the following is/are correct
\begin{qoptions}
\opt{A}{$x+2z=0$}
\opt{B}{$x=2y=2z=1$}
\opt{C}{$y+z=0$}
\opt{D}{$x=2y$}
\end{qoptions}
\qn{2} $\displaystyle\lim_{n\to\infty}\left(\sin^n 1^c+\cos^n 1^c\right)^{1/n}$, is
\begin{qoptions}
\opt{A}{greater than $\dfrac{1}{\sqrt2}$}
\opt{B}{greater than $\dfrac{\sqrt3}{2}$}
\opt{C}{less than $\dfrac{\sqrt3}{2}$}
\opt{D}{equal to 1}
\end{qoptions}
\qn{3} Let P denotes the plane consisting of all points that are equidistant from the points A$(-4,2,1)$ and B$(2,-4,3)$ and Q be the plane $x-y+cz=1$ where $c\in\mathbb{R}$.
\begin{qoptions}
\opt{A}{The plane P is parallel to plane Q if $c=1/3$}
\opt{B}{The plane P is parallel to plane Q if $c=3$}
\opt{C}{If angle between the planes P \& Q is $45^\circ$, then product of all possible values of `c' is $-2$}
\opt{D}{If angle between the planes P \& Q is $45^\circ$, then sum of all possible values of `c' is $\dfrac{24}{17}$}
\end{qoptions}
\qn{4} If $\displaystyle\lim_{x\to0}\dfrac{ae^{1/x}+be^{-1/x}}{ce^{1/x}+de^{-1/x}}=2$, then the roots of $bx^2+(a-2c)x-2d=0$; $a,b,c,d\in\mathbb{Q}$ are
\begin{qoptions}
\opt{A}{real}
\opt{B}{integral}
\opt{C}{rational}
\opt{D}{equal}
\end{qoptions}
\qn{5} If $f(x)$ is a continuous function for all $x\in\mathbb{R}$ and the range of $f(x)$ is $(2,\sqrt{26})$ and $g(x)=\left[\dfrac{f(x)}{a}\right]$ is continuous for all $x\in\mathbb{R}$ ([.] denotes the greatest integer function), then which of the following cannot be a value of `a' :
\begin{qoptions}
\opt{A}{2}
\opt{B}{3}
\opt{C}{6}
\opt{D}{5}
\end{qoptions}
\qn{6} If $f(x)=\cos^{-1}\left(\dfrac{x^{2n}-1}{x^{2n}+1}\right)$, then $f'(x)=$
\begin{qoptions}
\opt{A}{$\dfrac{-2nx^{n-1}}{x^{2n}+1}$, if $n$ is even}
\opt{B}{$\dfrac{-2nx^{n}}{|x|\left(x^{2n}+1\right)}$, if $n$ odd \& $x\neq0$}
\opt{C}{$\dfrac{-2nx^{n}}{|x|\left(x^{2n}+1\right)}$, if $n$ even \& $x\neq0$}
\opt{D}{$\dfrac{-2nx^{n-1}}{x^{2n}+1}$, if $n$ is odd}
\end{qoptions}
\columnbreak
\sect{Section-I(ii)}
\qn{7}\par\nopagebreak[4]
\begin{center}\small\setlength{\tabcolsep}{3pt}
\begin{tabular}{|c|p{4.2cm}|c|c|}
\hline
\multicolumn{2}{|c|}{List-I} & \multicolumn{2}{c|}{List-II}\\ \hline
(P) & If $\vec a$ and $\vec b$ are any two unit vectors, then possible integer in the range of $\dfrac{3\left|\vec a+\vec b\right|}{2}+2\left|\vec a-\vec b\right|$ can be & (1) & $0$\\ \hline
(Q) & If $\hat a,\hat b$ and $\hat c$ are non coplanar vectors then the vectors $\vec V_1=\hat a+2\hat b+3\hat c$; $\vec V_2=\lambda\hat b+4\hat c$ and $\vec V_3=(2\lambda-1)\hat c$ where $\lambda$ is a scalar, can be coplanar, for $\lambda$ equal to & (2) & $1$\\ \hline
(R) & The distance of the z-axis from the image of the point A$(2,-3,3)$ in the plane $x-2y-z+1=0$, is & (3) & $2$\\ \hline
(S) & The figure shows a pyramid OABCD (where O is the origin). The base OABC is a square of side 1 unit. The point D is at a distance of 1 unit from the square base. D lies on the normal passing through the centroid of the base. If the angle between $\overrightarrow{\mathrm{OB}}$ and $\overrightarrow{\mathrm{OD}}$ is $\tan^{-1}\sqrt{K}$, then K is equal to\par\smallskip\qfigi{p1_math_s2_q07_S} & (4) & $-1$\\ \hline
 & & (5) & $4$\\ \hline
\end{tabular}
\end{center}
\begin{qoptions}
\opt{A}{P $\to$ 4;Q $\to$ 2;R $\to$ 1;S $\to$ 5}
\opt{B}{P $\to$ 5;Q $\to$ 1;R $\to$ 2;S $\to$ 3}
\opt{C}{P $\to$ 5;Q $\to$ 1;R $\to$ 3;S $\to$ 2}
\opt{D}{P $\to$ 4;Q $\to$ 1;R $\to$ 2;S $\to$ 5}
\end{qoptions}
\columnbreak
\qn{8} Let $f_1(x)$ and $f_2(x)$ be two linear polynomials defined from $[-1,1]$ onto $[0,2]$. $f_1(x)$ is an increasing function and $f_2(x)$ is a decreasing function.\par\nopagebreak[4]
\begin{center}\small\setlength{\tabcolsep}{3pt}
\begin{tabular}{|c|p{3.9cm}|c|c|}
\hline
\multicolumn{2}{|c|}{List-I} & \multicolumn{2}{c|}{List-II}\\ \hline
(P) & If N is the number of dissimilar terms in $(f_1(x))^{2010}+(f_2(x))^{2010}$ then the value of `r' for which $\dbinom{N-1000}{r}$ is greatest, is (where $\dbinom{k}{r}$ equals ${}^k\mathrm{C}_r$) & (1) & $1$\\ \hline
(Q) & The term independent of `x' in the expansion of $\left(9x-\dfrac{1}{3\sqrt{x}}\right)^{18}$, $x>0$, is $\alpha$ times the corresponding binomial co-efficient. Then `$\alpha$' is & (2) & $2$\\ \hline
(R) & If $\displaystyle\left|\sum_{x=2}^{\infty}\dfrac{1}{f_1(x)\,f_2(x)}\right|=\dfrac{3}{n}$, then $n$ is & (3) & $3$\\ \hline
(S) & If $|f_1(z)|=1$ and $|f_2(z)|=2$, where $z$ is a complex number, then the number of values of $z$ consistent with both equations equals & (4) & $4$\\ \hline
 & & (5) & $5$\\ \hline
\end{tabular}
\end{center}
\begin{qoptions}
\opt{A}{P $\to$ 3;Q $\to$ 2;R $\to$ 4;S $\to$ 5}
\opt{B}{P $\to$ 5;Q $\to$ 1;R $\to$ 5;S $\to$ 3}
\opt{C}{P $\to$ 3;Q $\to$ 1;R $\to$ 4;S $\to$ 2}
\opt{D}{P $\to$ 4;Q $\to$ 3;R $\to$ 1;S $\to$ 2}
\end{qoptions}
\qn{9}\par\nopagebreak[4]
\begin{center}\small\setlength{\tabcolsep}{3pt}
\begin{tabular}{|c|c|c|l|}
\hline
\multicolumn{2}{|c|}{List-I} & \multicolumn{2}{c|}{List-II}\\ \hline
(P) & $\cot^{-1}\left(\tan(-37^\circ)\right)$ & (1) & $143^\circ$\\ \hline
(Q) & $\cos^{-1}\left(\cos(-233^\circ)\right)$ & (2) & $127^\circ$\\ \hline
(R) & $\sin\left(\dfrac12\cos^{-1}\left(\dfrac19\right)\right)$ & (3) & $\dfrac34$\\ \hline
(S) & $\cos\left(\dfrac12\operatorname{arc\,cos}\left(\dfrac18\right)\right)$ & (4) & $\dfrac23$\\ \hline
 & & (5) & None of the above\\ \hline
\end{tabular}
\end{center}
\begin{qoptions}
\opt{A}{P $\to$ 2;Q $\to$ 2;R $\to$ 5;S $\to$ 4}
\opt{B}{P $\to$ 1;Q $\to$ 5;R $\to$ 4;S $\to$ 5}
\opt{C}{P $\to$ 5;Q $\to$ 1;R $\to$ 3;S $\to$ 3}
\opt{D}{P $\to$ 2;Q $\to$ 2;R $\to$ 4;S $\to$ 3}
\end{qoptions}
\columnbreak
\qn{10}\par\nopagebreak[4]
\begin{center}\small\setlength{\tabcolsep}{3pt}
\begin{tabular}{|c|p{4.2cm}|c|c|}
\hline
\multicolumn{2}{|c|}{List-I} & \multicolumn{2}{c|}{List-II}\\ \hline
(P) & If $\cot^{-1}x+\cot^{-1}y+\cot^{-1}z=\dfrac{\pi}{4}$, then $xy+yz+zx+x+y+z-xyz=$ & (1) & $3$\\ \hline
(Q) & The number of solution(s) of the equation $2\left(\sin^{-1}x\right)^2-5\sin^{-1}x+2=0$ is & (2) & $0$\\ \hline
(R) & If $\sin^{-1}2x-\cos^{-1}2x=\sin^{-1}(6x-2)$, then number of values of $x$ is & (3) & $2$\\ \hline
(S) & If $3\tan^{-1}(2-\sqrt3)-\tan^{-1}\dfrac1x=\tan^{-1}\left(\dfrac13\right)$, then $x=$ & (4) & $1$\\ \hline
\end{tabular}
\end{center}
\begin{qoptions}
\opt{A}{P $\to$ 4;Q $\to$ 4;R $\to$ 3;S $\to$ 3}
\opt{B}{P $\to$ 4;Q $\to$ 4;R $\to$ 1;S $\to$ 3}
\opt{C}{P $\to$ 1;Q $\to$ 2;R $\to$ 4;S $\to$ 3}
\opt{D}{P $\to$ 2;Q $\to$ 1;R $\to$ 3;S $\to$ 3,4}
\end{qoptions}
\sect{Section-II}
\qn{1} Consider two lines $L_1:\dfrac{x-7}{3}=\dfrac{y-7}{2}=\dfrac{z-3}{1}$ and $L_2:\dfrac{x-1}{2}=\dfrac{y+1}{4}=\dfrac{z+1}{3}$. If a line L whose direction ratios are $<2,2,1>$ intersect the line $L_1$ and $L_2$ at A and B, then find the distance AB.
\qn{2} A twice differentiable function $f(x)$ is defined for all real numbers and satisfies the following conditions $f(0)=2$; $f'(0)=-5$ and $f''(0)=3$. The function $g(x)$ is defined by $g(x)=e^{ax}+f(x)\ \ \forall\ x\in\mathbb{R}$, where `a' is a positive constant. If $g'(0)+g''(0)=0$. Find the value of `a'.
\qn{3} Suppose that $f$ and $g$ are non--constant continuous and differentiable, real valued functions on $\mathbb{R}$. If for every $x,y\in\mathbb{R}$, $f(x+y)=f(x).f(y)-g(x).g(y)$ and $g(x+y)=g(x).f(y)+f(x).g(y)$. Also $f'(0)=0$ then find the value of $\dfrac{1-\left\{\left(f(x)\right)^6+\left(g(x)\right)^6\right\}}{\left(f(x)\right)^2\left(g(x)\right)^2}$.
\qn{4} Let $A_1,A_2,A_3,A_4$ be the areas of the triangular faces of a tetrahedron and $h_1,h_2,h_3,h_4$ be the corresponding altitudes drawn from the opposite vertices of the tetrahedron. If volume of tetrahedron is 5 cubic units, then find the minimum value of $(A_1+A_2+A_3+A_4)(h_1+h_2+h_3+h_4)$ (in cubic units).
\qn{5} If $0<\cos^{-1}x<1$ and $1+\sin(\cos^{-1}x)+\sin^2(\cos^{-1}x)+\sin^3(\cos^{-1}x)+\cdots\infty=2$, then find the value of $x^2$.
\qn{6} A student forgot the product rule for differentiation and made the mistake of thinking that $(f\cdot g)'=f'g'$. However he was lucky to get the correct answer. The function $f$ that he used was $f(x)=e^{x^2}$. If the domain of $g(x)$ is the interval $\left(\dfrac12,\infty\right)$ with $g(1)=e$, then the value of $g(5)e^{-5}$ is
\qn{7} Let ABCD be any quadrilateral. P and Q are the midpoints of its diagonal. If $\left|\overrightarrow{AB}\right|^2+\left|\overrightarrow{BC}\right|^2+\left|\overrightarrow{CD}\right|^2+\left|\overrightarrow{DA}\right|^2-\left|\overrightarrow{AC}\right|^2-\left|\overrightarrow{BD}\right|^2=\lambda\left|\overrightarrow{PQ}\right|^2$, then find the value of $\lambda$.
\qn{8} A sequence $x_1,x_2,x_3,\ldots$ of positive real numbers is such that $x_1=\tan^{-1}2$ and $\sin(x_{n+1}-x_n)+2^{-(n+1)}\sin x_n\cdot\sin x_{n+1}=0$ for all $n\geqslant1$ further $x_1>x_2>x_3>\ldots>x_{n+1}>\ldots$. Find the greatest integral value of `n', such that $\cot x_n<\dfrac{1023}{1024}$.
\end{multicols}
\end{document}
